\section*{Introduction}
Artificial intelligence (AI) has made significant strides in recent decades, promising to impact myriad aspects of society, including education~\supercite{adiguzel2023revolutionizing,akgun2022artificial}, healthcare~\supercite{mesko2023imperative,loh2022application}, and industry~\supercite{ahmed2022artificial}.
Major investments in predictive and generative AI have catalyzed society-level debates over the future of AI at home and in the workplace.
Perhaps more than any other domain, AI tools have become deeply entwined with the process of knowledge production, yielding findings that attract disproportionate attention in various scientific fields~\supercite{wang2023scientific}.
For example, AlphaFold learns known protein structures to accurately predict the unexplored ones, circumventing the capital and human cost of conventional structural inference and recently granted a 2024 Nobel Prize~\supercite{jumper2021highly,varadi2022alphafold}.
Models improved via deep reinforcement learning have become tuned to contain complex fusion reactions~\supercite{degrave2022magnetic} and discovered new, hardware-optimized forms to matrix multiplication that recursively accelerate deep learning itself~\supercite{fawzi2022discovering}. 
Autonomous laboratory systems driven by ChatGPT have helped some chemists and material scientists upscale the number of adaptive high-throughput experiments~\supercite{boiko2023autonomous,stokel2023chatgpt,gilson2023does}.
Moreover, recent developments in large language models are making them increasingly incorporated in assisting scientific writing~\supercite{salimi2023large,liang2024mapping,hwang2024can,kobak2025delving}, facilitating the distillation of scientific findings, but also raising concerns about weakened confidence in AI-generated content~\supercite{stokel2023chatgpt,gilson2023does, wojtowicz2025undermining}.
AI's increasing capabilities to influence scientific research suggest that it manifests potential to both increase the productivity of individual scientists and raise the visibility of science it supports.

Despite the increasing adoption of AI in science, large-scale empirical measurements of AI’s scientific impact are limited, and a detailed, dynamic understanding of AI’s impact on the entire character of science remains largely unknown. Recent work suggests AI has brought widespread benefits to individual scientists but may lead to demographic disparity resulting from gaps in AI education~\supercite{gao2024quantifying}. Researchers have also identified evolving citation patterns that signal a changing scientific landscape in AI research~\supercite{frank2019evolution}.
Here we seek to explore the impact of AI in scientific research at different scales by posing the question: How does the adoption of AI influence individual scientists’ careers and the collective exploration of science as a whole?

We conduct a large-scale quantitative analysis of the impact of AI on scientists and science, covering 41,298,433 research papers spanning from 1980 to 2025 in the OpenAlex dataset~\supercite{openalex}, with patterns corroborated using the Web of Science~\supercite{WoS,mongeon2016journal}.
Notably, we do not focus on computer science or mathematics, fields that develop AI methodologies directly, but rather on papers that augment research in natural science fields by adopting AI, primarily covering decades involving development and deployment of conventional machine learning algorithms and also extending to a necessarily more preliminary analysis of the latest generative AI techniques.
Specifically, we select six representative disciplines that cover the vast majority of natural science contributions—biology, medicine, chemistry, physics, materials science, and geology.
We then leverage the BERT language model~\supercite{DBLP:conf/naacl/DevlinCLT19,wolf-etal-2020-transformers} to accurately identify such AI-augmented research papers based on their titles and abstracts.

We separate the periods in which AI was predominantly conventional machine learning, deep learning, and most recently generative designs including large language models. With abundant data-based evidence across decades of conventional machine learning and deep learning, we validate these AI-based measurements and use them to reveal that the adoption of AI leads to an amplifying effect on the career of individual scientists, bringing acceleration in the production and visibility of science produced by those scientists who incorporate AI.
Nevertheless, this effect corresponds with a contracted focus within collective science.
Measured with “knowledge extent”, the “diameter” covered by a sampled batch of papers in vector space, AI-driven science spans less topical ground and is associated with a decrease in follow-on scientific engagement, suggesting that AI is currently more likely to focus on existing popular research problems rather than explore new ones.
Meanwhile, analyses using currently available data within the latest era of generative AI including large language models reveal a preliminary consistency with prior periods, providing a starting point for further study as generative AI develops over a longer period.